\chapter*{第 3 章 流体运动学}
\addcontentsline{toc}{chapter}{第 3 章 流体运动学}
\markboth{第 3 章 流体运动学}{}

\begin{tishi}
本章是 818 全卷\textbf{出题面最宽}的一章：真题不在本章出计算大题，但选择、填空、判断、简答的
名词定义与概念辨析大量取自本章，另有一类小计算（流线方程、加速度分量、连续性方程求流速与密度）。
本册按教材原顺序逐题排列：\textbf{第 1 题＝习题 3.1，第 2 题＝习题 3.2，$\cdots$，第 11 题＝习题 3.11}。
做题要求：概念题\textbf{写出完整要点}（不是“略”），计算题必须写出公式、代入、单位、结果四步。
\end{tishi}

\begin{ti}
\sloppy
\emergencystretch=2em

\item 选择题（共 4 小题）。

\begin{xiaowen}
\item 恒定流（定常流）是（\quad）。

\begin{choices}
\item 流动随时间按一定规律变化
\item 流场中各空间点的运动参数不随时间变化
\item 各过流断面的速度分布相同
\item 各过流断面的压强相同
\end{choices}

\item 非恒定流是（\quad）。

\begin{choices}
\item $\dfrac{\partial u}{\partial t}=0$
\item $\dfrac{\partial u}{\partial t}\neq0$
\item $\dfrac{\partial u}{\partial s}=0$
\item $\dfrac{\partial u}{\partial s}\neq0$
\end{choices}

\item 渐变流是（\quad）。

\begin{choices}
\item 过流断面上速度均匀分布的流动
\item 速度不随时间变化的流动
\item 流线近于平行直线的流动
\item 沿程各断面流量相同的流动
\end{choices}

\item 动能修正系数是反映过流断面上实际流速分布不均匀性的系数，流速分布（\quad），
系数值（\quad）；流速分布（\quad）时，则动能修正系数接近于（\quad）。

\begin{choices}
\item 越不均匀，越小；均匀，1
\item 越均匀，越小；均匀，1
\item 越不均匀，越小；均匀，0
\item 越均匀，越小；均匀，0
\end{choices}
\end{xiaowen}
\xstarfull{3}\yiju{E1 slide16「选择/判断考概念辨析；填空考名词定义与小计算」；E2「第 1--3 章偏概念，掌握基本公式与解释即可」；（4）动能修正系数系第 4 章伯努利方程的修正项、真题中未单独出现，其重要度按两星处理}\nandu{易}

\item 用欧拉观点写出下列各情况下密度变化率的数学表达式：（1）均质流体；（2）不可压缩流体；（3）不可压缩均质流体。
\xstarfull{3}\yiju{E1 slide33「第三章流体运动学＝公式多、重点记忆」；E2「第 1--3 章偏概念，掌握基本公式与解释即可」；教材式 3.9--3.11 与密度随体导数的展开式（E6）}\nandu{中}

\item 已知三维速度场 $u_x=yz+t$，$u_y=xz-t$，$u_z=xy$。（1）判断流动是否恒定；（2）求流体质点在通过流场点 $(1,1,1)$ 时的加速度。
\xstarfull{4}\yiju{E1 slide33「第三章＝公式多、重点记忆」；E2「第 1--3 章……掌握基本公式」；质点加速度式（教材式 3.7、3.8）是第 4 章欧拉运动微分方程与伯努利方程的直接输入（E6）}\nandu{中}

\item 已知二维速度场 $u_x=x+2t$，$u_y=-y+t-3$。试求：该流动的流线方程以及在 $t=0$ 瞬时过点 $M(-1,-1)$ 的流线。
\xstarfull{4}\yiju{E4 2015 年试题计算题第 1 题（流线方程与流动方向判断）；E1 slide33「第三章＝公式多、重点记忆」；教材式 3.14（E6）}\nandu{中}

\item 已知二维非恒定流场的速度分布为：$u_x=x+t$，$u_y=-y+t$。试求：$t=0$ 和 $t=2$ 时，过点 $M(-1,-1)$ 的流线方程。
\xstarfull{4}\yiju{E4 2015 年试题计算题第 1 题（流线方程，且为非恒定流场）；E1 slide33「第三章＝公式多、重点记忆」；教材式 3.14（E6）}\nandu{中}

\item 如习题 3.6 图所示，已知两平板间的速度分布为 $u_x=u_{\max}\left[1-\left(\dfrac{y}{b}\right)^2\right]$，式中，$y=0$ 为中心线，$y=\pm b$ 为平板所在位置，$u_{\max}$ 为常数，试求通过平板的流体单宽流量。
\tuyi{习题 3.6 图：上、下两块固定平板，板面平行于 $x$ 轴、间距为 $2b$；$y$ 轴垂直板面，原点 $O$ 在中心线上，故上板位于 $y=+b$、下板位于 $y=-b$；两板间流体沿 $x$ 方向流动，断面流速呈抛物线分布，顶点在中心线 $y=0$ 处、最大值为 $u_{\max}$，在两板处降为零。所求单宽流量即沿 $y$ 从 $-b$ 积到 $+b$ 的流量。}
\xstarfull{2}\yiju{E2「第 1--3 章偏概念」；教材式 3.15、3.16（流量积分式，E6）；E4 历年真题中未出现同类积分求流量题，故压至两星}\nandu{中}

\item 试证以下不可压缩流体的运动是可能存在的：$u_x=2x^2+y$，$u_y=2y^2+z$，$u_z=-4(x+y)z+xy$。
\xstarfull{3}\yiju{E4 2022 年 818 单选第 4 题（连续性方程的同型应用）；E1 slide33「第三章＝公式多、重点记忆」；教材式 3.25（E6）}\nandu{中}

\item 已知平面不可压缩流体速度分布：$u_x=2x-1$，$u_y=-2y$。（1）流动满足连续性方程否？（2）流动是否无旋？（3）求加速度、线变形率和纯剪切变形率。
\xstarfull{4}\yiju{E5 题库 \texttt{10.pdf} 选择第 13 题（「涡量为零」的有旋无旋判别）；E4 2022 年 818 单选第 4 题（连续性方程）；E1 slide33「第三章＝公式多、重点记忆」}\nandu{中}

\item 变直径管的直径 $d_1=320\ \mathrm{mm}$，$d_2=160\ \mathrm{mm}$，断面平均流速 $v_1=1.5\ \mathrm{m/s}$，那么断面平均流速 $v_2$ 应等于多少？
\xstarfull{3}\yiju{E4 2015 年试题填空（变径管求 $v_2$）；E2「第 1--3 章……掌握基本公式」；教材式 3.31（E6）}\nandu{易}

\item 空气从断面积 $A_1=0.4\ \mathrm{m}\times0.4\ \mathrm{m}$ 的方形管中进入压缩机，密度 $\rho_1=1.2\ \mathrm{kg/m^3}$，断面平均流速 $v_1=4\ \mathrm{m/s}$。压缩后，从直径 $d_1=0.25\ \mathrm{m}$ 的圆形管中排出，断面平均流速 $v_2=3\ \mathrm{m/s}$。试求：压缩机出口断面的平均密度 $\rho_2$ 和质量流量 $Q_m$。
\xstarfull{5}\yiju{E4 2022 年 818 单选第 4 题（方管进口接压缩机，求出口平均密度 $\rho_2$ 与质量流量 $Q_m$）；E1 slide33「第三章＝公式多、重点记忆」；教材式 3.30（E6）}\nandu{中}

\item 如习题 3.11 图所示，输水管道经三通管汇流，已知流量 $Q_1=1.5\ \mathrm{m^3/s}$，$Q_3=2.6\ \mathrm{m^3/s}$，过流断面面积 $A_2=0.2\ \mathrm{m^2}$，试断面平均流速 $v_2$。
\tuyi{习题 3.11 图：左侧两根支管斜向汇入一根水平干管。断面 1--1 位于上支管上、流量为 $Q_1$；断面 2--2 位于下支管上、流量为 $Q_2$、过流断面面积为 $A_2$；两支管汇合后由右侧水平干管的断面 3--3 流出，流量为 $Q_3$。}
\xstarfull{3}\yiju{E2「第 1--3 章……掌握基本公式与解释即可」；教材式 3.32（分流与汇流的总流连续性方程，E6）；E4 历年真题中未直接出现三通汇流题，故按三星处理}\nandu{易}

\end{ti}

\begin{banfa}
\textbf{按题号给出的应试提示}\\
\textbf{第 3.2 题：}欧拉观点下“某物理量的变化率”就是它的随体导数，先写出
$\dfrac{\dd\rho}{\dd t}=\dfrac{\partial\rho}{\partial t}+(\bu\cdot\grad)\rho$，再按“与空间无关”“随体不变”“二者兼有”逐条删项。\\
\textbf{第 3.3、3.8 题：}求加速度三步：写出三个速度分量 $\to$ 算出九个一阶偏导 $\to$ 按
$a_x=\dfrac{\partial u_x}{\partial t}+u_x\dfrac{\partial u_x}{\partial x}+u_y\dfrac{\partial u_x}{\partial y}+u_z\dfrac{\partial u_x}{\partial z}$
代入；定常时先令 $\dfrac{\partial(\cdot)}{\partial t}=0$，只剩三项迁移加速度。\\
\textbf{第 3.4、3.5 题：}流线用 $\dfrac{\dd x}{u_x}=\dfrac{\dd y}{u_y}$，积分时 $t$ 当常数；
问到“某瞬时过某点的流线”，先把该瞬时代入消去 $t$，再用定点坐标定出积分常数。\\
\textbf{第 3.6 题：}单宽流量 $q=\displaystyle\int_{-b}^{b}u_x\dd y$，对称区间直接积，结果 $\dfrac{4}{3}u_{\max}b$；
判别断面平均流速时用 $v=\dfrac{q}{2b}=\dfrac{2}{3}u_{\max}$（不是 $\dfrac{u_{\max}}{2}$）。\\
\textbf{第 3.7、3.8 题：}连续性只验 $\grad\cdot\bu$：恒为零即“运动可能存在”；再判断无旋就看
$\omega_x,\omega_y,\omega_z$ 是否全为零。\\
\textbf{第 3.9--3.11 题：}一元总流连续性——不可压缩 $v_1A_1=v_2A_2$（流速与面积成反比、与直径平方成反比）；
气体 ${\rho_1v_1A_1=\rho_2v_2A_2}$；有汇流时写成 $Q_1+Q_2=Q_3$（“进＝出”）。
\end{banfa}

\begin{yicuo}
\textbf{本章三个最高频的陷阱：}\\
① \textbf{恒定流与均匀流是两回事。}恒定流看\textbf{时间}（$\dfrac{\partial(\cdot)}{\partial t}=0$），
均匀流看\textbf{空间}（流线是相互平行的直线）；恒定流可以是非均匀流，非恒定流也可以是均匀流。\\
② \textbf{求流线时 $t$ 是常数，求迹线时 $t$ 是自变量。}把 $t$ 当自变量积进去得到的是迹线，不是流线。\\
③ \textbf{不可压缩与均质（$\rho=$ 常数）不是一回事。}不可压缩写 $\dfrac{\dd\rho}{\dd t}=0$，
只说明\textbf{同一质点}在运动全过程中密度不变，\textbf{不同质点可以有不同的密度}，
所以不可压缩流体不一定处处等密度。
\end{yicuo}
